\Needspace{8\baselineskip}
\Article{208}
\textsc{Aufgabe.} \textit{Für zwei gegebene eigentlich äquivalente Formen $F$, $F'$ mit der Determinante $hh$ eine eigentliche Transformation der einen in die andere finden.}

\SourceLeaf{207}
\textit{Lösung.} Zur Form $F$ sei die reduzierte Form $\Phi$ eigentlich äquivalent, die somit nach Voraussetzung auch zu $F'$ eigentlich äquivalent ist. Wir suchen nach Artikel~206 eine eigentliche Transformation von $F$ in $\Phi$, etwa $\alpha,\beta,\gamma,\delta$, und ebenso eine eigentliche Transformation von $F'$ in $\Phi$, etwa $\alpha',\beta',\gamma',\delta'$. Dann geht $\Phi$ in $F'$ über, und zwar durch die eigentliche Substitution $\delta',-\beta',-\gamma',\alpha'$, und somit $F$ in $F'$ durch die eigentliche Substitution
\[
\alpha\delta'-\beta\gamma',\quad\beta\alpha'-\alpha\beta',\quad
\gamma\delta'-\delta\gamma',\quad\delta\alpha'-\gamma\beta'
\]

Es lohnt sich, für diese Transformation von $F$ in $F'$ eine andere Formel herzuleiten, für die wir nicht einmal die reduzierte Form $\Phi$ selbst kennen müssen. Wir nehmen an:
\[
F\text{ sei }(a,b,c),\quad F'=(a',b',c'),\quad\Phi=(A,h,0)
\]
Da den Verhältnissen $h-b:a$ oder $c:-(h+b)$ in kleinsten Zahlen das Verhältnis $\beta:\delta$ gleich ist, sehen wir leicht, dass $\frac{h-b}{\beta}=\frac a\delta$ eine \textit{ganze Zahl} ist; sie sei $f$. Ebenso ist $\frac c\beta=\frac{-h-b}{\delta}$ eine ganze Zahl; wir setzen sie $=g$. Wir haben
\[
\begin{gathered}A=a\alpha\alpha+2b\alpha\gamma+c\gamma\gamma\quad\text{also}\\
\beta A=a\alpha\alpha\beta+2b\alpha\beta\gamma+c\beta\gamma\gamma\end{gathered}
\]
oder, wenn wir für $a\beta$ den Wert $\delta(h-b)$ und für $c$ den Wert $\beta g$ einsetzen,
\[
\beta A=\alpha\alpha\delta h+b(2\beta\gamma-\alpha\delta)\alpha+\beta\beta\gamma\gamma g
\]
oder wegen $b=-h-\delta g$
\[
\beta A=2\alpha(\alpha\delta-\beta\gamma)h+(\alpha\delta-\beta\gamma)^2g=2\alpha h+g
\]
Ebenso gilt
\[
\begin{aligned}
\delta A&=a\alpha\alpha\delta+2b\alpha\gamma\delta+c\gamma\gamma\delta\\
&=\alpha\alpha\delta\delta f+b(2\alpha\delta-\beta\gamma)\gamma-\beta\gamma\gamma h\\
&=(\alpha\delta-\beta\gamma)^2f+2\gamma(\alpha\delta-\beta\gamma)h=2\gamma h+f
\end{aligned}
\]
Also
\[
\alpha=\tfrac{\beta A-g}{2h},\qquad\gamma=\tfrac{\delta A-f}{2h}
\]
Setzen wir auf genau dieselbe Weise
\[
\frac{h-b'}{\beta'}=\frac{a'}{\delta'}=f',\qquad
\frac{c'}{\beta'}=\frac{-h-b'}{\delta'}=g'
\]
so ergibt sich
\[
\alpha'=\tfrac{\beta' A-g'}{2h},\qquad\gamma'=\tfrac{\delta' A-f'}{2h}
\]
\SourceLeaf{208}
Setzen wir diese Werte von $\alpha,\gamma,\alpha',\gamma'$ in die eben angegebene Formel für die Transformation von $F$ in $F'$ ein, so geht sie über in
\[
\tfrac{\beta f'-\delta'g}{2h},\quad
\tfrac{\beta'g-\beta g'}{2h},\quad
\tfrac{\delta f'-\delta'f}{2h},\quad
\tfrac{\beta'f-\delta g'}{2h}
\]
woraus $A$ ganz verschwunden ist.

Sind zwei uneigentlich äquivalente Formen $F$, $F'$ gegeben und suchen wir eine uneigentliche Transformation der einen in die andere, so sei $G$ die zu $F$ entgegengesetzte Form und eine eigentliche Transformation von $G$ in $F'$ die folgende: $\alpha,\beta,\gamma.\ \delta$. Dann ist offenbar $\alpha,\beta,-\gamma,-\delta$ eine uneigentliche Transformation von $F$ in $F'$.

Sind schließlich die gegebenen Formen sowohl eigentlich als auch uneigentlich äquivalent, so können wir auf diese Weise offenbar zwei Transformationen finden, eine eigentliche und eine uneigentliche.

\Needspace{8\baselineskip}
\Article{209}
Es bleibt also nur noch zu zeigen, wie sich aus einer Transformation alle übrigen gleichartigen ableiten lassen. Das hängt von der Lösung der unbestimmten Gleichung $tt-hhuu=mm$ ab, wobei $m$ den größten gemeinsamen Teiler von $a,2b,c$ bezeichnet, wenn $(a,b,c)$ eine der beiden äquivalenten Formen ist. Diese Gleichung lässt sich aber stets nur auf zwei Arten lösen: Wir setzen entweder $t=m$, $u=0$ oder $t=-m$, $u=0$. Nehmen wir nämlich an, es gebe noch eine andere Lösung $t=T$, $u=U$, wobei $U$ nicht $=0$ ist. Da $mm$ sicher $4hh$ teilt, ist $\frac{4TT}{mm}=\frac{4hhUU}{mm}+4$, und sowohl $\frac{4TT}{mm}$ als auch $\frac{4hhUU}{mm}$ sind ganze Quadratzahlen. Nun sehen wir leicht, dass $4$ nur dann die Differenz zweier ganzer Quadratzahlen sein kann, wenn das kleinere Quadrat $0$ ist, \textit{d.h.} $U=0$, entgegen der Voraussetzung. --- Geht also die Form $F$ in $F'$ über, und zwar durch die Substitution $\alpha,\beta,\gamma,\delta$, so gibt es außer $-\alpha,-\beta,-\gamma,-\delta$ keine weitere dazu gleichartige Transformation. Sind daher zwei Formen entweder nur eigentlich oder nur uneigentlich äquivalent, so gibt es nur \textit{zwei} Transformationen; sind sie sowohl eigentlich als auch uneigentlich äquivalent, so gibt es \textit{vier}, nämlich zwei eigentliche und zwei uneigentliche.

\Needspace{8\baselineskip}
\Article{210}
\textsc{Satz.} \textit{Sind zwei reduzierte Formen $(a,h,0)$, $(a',h,0)$ uneigentlich äquivalent, so gilt $aa'\equiv mm\;(\mathit{mod.}\,2mh)$, wobei $m$ den größten gemeinsamen Teiler von $a,2h$ oder $a',2h$ bezeichnet. Haben umgekehrt $a,2h;\ a',2h$ denselben größten gemeinsamen}
\SourceLeaf{209}
\textit{Teiler $m$ und gilt $aa'\equiv mm\;(\mathit{mod.}\,2mh)$, so sind die Formen $(a,h,0)$, $(a',h,0)$ uneigentlich äquivalent.}

\textit{Beweis.} I. Die Form $(a,h,0)$ gehe in $(a',h,0)$ über, und zwar durch die uneigentliche Substitution $\alpha,\beta,\gamma,\delta$, sodass wir vier Gleichungen haben:
\begin{align*}
a\alpha\alpha+2h\alpha\gamma&=a'\tag*{. . . . . . . . . . . . [1]}\\
a\alpha\beta+h(\alpha\delta+\beta\gamma)&=h\tag*{. . . . . . . . . . [2]}\\
a\beta\beta+2h\beta\delta&=0\tag*{. . . . . . . . . . . . [3]}\\
\alpha\delta-\beta\gamma&=-1\tag*{. . . . . . . . . . . . . [4]}
\end{align*}
Multiplizieren wir [4] mit $h$ und ziehen es von [2] ab, was wir durch $[2]-h[4]$ ausdrücken, so ergibt sich
\begin{equation*}
(a\alpha+2h\gamma)\beta=2h\tag*{. . . . . . . . . . . . [5]}
\end{equation*}
Ebenso folgt aus $\gamma\delta[2]-\gamma\gamma[3]-(a+a\beta\gamma+h\gamma\delta)[4]$, nachdem die sich aufhebenden Glieder wegfallen,
\begin{equation*}
-a\alpha\delta=a+2h\gamma\delta,\quad\text{oder}\quad -(a\alpha+2h\gamma)\delta=a\tag*{. . [6]}
\end{equation*}
Schließlich erhalten wir $a[1]\ldots a\alpha(a\alpha+2h\gamma)=aa'$, oder
\[
(a\alpha+2h\gamma)^2-aa'=2h\gamma(a\alpha+2h\gamma)
\]
oder
\begin{equation*}
(a\alpha+2h\gamma)^2\equiv aa'\;(\mathrm{mod.}\,2h(a\alpha+2h\gamma))\tag*{. . . . . [7]}
\end{equation*}
Aus [5] und [6] folgt nun, dass $a\alpha+2h\gamma$ sowohl $2h$ als auch $a$ und somit auch $m$ teilt, den größten gemeinsamen Teiler von $a,2h$. Offenbar teilt aber $m$ auch $a\alpha+2h\gamma$. Also muss $a\alpha+2h\gamma$ entweder $=+m$ oder $=-m$ sein. Daraus folgt nach [7] sofort $mm\equiv aa'\;(\mathrm{mod.}\,2mh)$. \textit{Damit ist die Behauptung bewiesen.}

II. Haben $a,2h;\ a',2h$ denselben größten gemeinsamen Teiler $m$ und ist außerdem $aa'\equiv mm\;(\mathrm{mod.}\,2mh)$, so sind $\frac a m$, $\frac{2h}{m}$, $\frac{a'}{m}$, $\frac{aa'-mm}{2mh}$ ganze Zahlen. Wir können leicht nachprüfen, dass die Form $(a,h,0)$ in $(a',h,0)$ übergeht, und zwar durch die Substitution $-\frac{a'}{m},-\frac{2h}{m},\frac{aa'-mm}{2mh},\frac a m$, und dass diese Transformation uneigentlich ist. Daher sind jene Formen uneigentlich äquivalent. \textit{Damit ist die Behauptung bewiesen.}

Damit lässt sich auch sofort entscheiden, ob eine gegebene reduzierte Form $(a,h,0)$ zu sich selbst uneigentlich äquivalent ist. Bezeichnen wir den größten gemeinsamen Teiler von $a,2h$ mit $m$, so muss $aa\equiv mm\;(\mathrm{mod.}\,2mh)$ gelten.

\SourceLeaf{210}
\Needspace{8\baselineskip}
\Article{211}
Alle reduzierten Formen mit gegebener Determinante $hh$ erhalten wir, indem wir in der unbestimmten Form $(A,h,0)$ für $A$ alle Zahlen von $0$ bis $2h-1$ einschließlich einsetzen. Ihre Anzahl ist also $2h$. Offenbar lassen sich alle Formen mit der Determinante $hh$ in ebenso viele \textit{Klassen} einteilen. Diese haben dieselben Eigenschaften, die wir oben (Artikel~175, 195) für die Klassen der Formen mit negativer und mit positiver nichtquadratischer Determinante besprochen haben. So zerfallen alle Formen mit der Determinante $25$ in zehn Klassen, die wir anhand der in jeder enthaltenen reduzierten Form unterscheiden können. Diese reduzierten Formen sind $(0,5,0)$, $(1,5,0)$, $(2,5,0)$, $(5,5,0)$, $(8,5,0)$, $(9,5,0)$, die zugleich zu sich selbst uneigentlich äquivalent sind; $(3,5,0)$, zu der $(7,5,0)$ uneigentlich äquivalent ist; und $(4,5,0)$, zu der $(6,5,0)$ uneigentlich äquivalent ist.

